% GENERATED FILE. Edit canonical lessons, then run npm run generate:books.
% canonical-source-hash: fnv1a64:2e0fb0d9346a9fcb
% canonical-lessons: PA-C210-perfective-return, PA-C210-ne, PA-C210-ne-review, PA-C210-aman-did-work, PA-C210-payoff

\chapter{Aman Did the Work}
\label{ch:pa-a2-ne}
\hlchaptermodality{210}

\begin{chapteropening}
\textbf{By the end of this chapter:} \emph{I can recognise and use the agent marker in one named-person completed-action frame.}

Recognise and produce one named-agent completed-action frame through listening, speaking, reading, and writing.
\end{chapteropening}

\section[\texorpdfstring{\pa{ਅਮਨ} · \pa{ਕੰਮ} · \pa{ਕੀਤਾ} (Aman; kamm; kiitaa)}{ਅਮਨ · ਕੰਮ · ਕੀਤਾ (Aman; kamm; kiitaa)}]{Return to a familiar completed action}

\label{lesson:PA-C210-perfective-return}



\begin{quote}
\emph{Hear:} \emph{Aman; kamm; kiitaa} Say \textbf{\pa{ਅਮਨ} · \pa{ਕੰਮ} · \pa{ਕੀਤਾ}}.
\end{quote}

\begin{grammarlens}[title={Grammar Lens}]
\textbf{\pa{ਕੰਮ}} is familiar “work.” \textbf{\pa{ਕੀਤਾ}} is the familiar completed form in \textbf{\pa{ਮੈਂ} \pa{ਕੰਮ} \pa{ਕੀਤਾ}}. In that old frame it agrees with the unmarked work object. Keep those pieces ready; the next lesson adds one small marker to a named person's completed action.
\end{grammarlens}

\subsection*{Your turn}
\begin{itemize}
\raggedright
  \item \emph{Hear:} \emph{the practice name} Say \textbf{\pa{ਅਮਨ}}.
  \item \emph{Hear:} \emph{work} Say \textbf{\pa{ਕੰਮ}}.
  \item \emph{Hear:} \emph{did, completed form} Say \textbf{\pa{ਕੀਤਾ}}.
\end{itemize}

\begin{tcolorbox}[breakable,colback=teal!4,colframe=teal!35!black,title={Before you move on}]
Say only these familiar pieces. No new form is required yet.
\end{tcolorbox}
\section[\texorpdfstring{\pa{ਨੇ} (ne)}{ਨੇ (ne)}]{A small marker after the agent}

\label{lesson:PA-C210-ne}



\begin{quote}
Recall the practice name \textbf{\pa{ਅਮਨ}}, “work” \textbf{\pa{ਕੰਮ}}, and the completed form \textbf{\pa{ਕੀਤਾ}}.
\end{quote}

\begin{grammarlens}[title={Grammar Lens: direct owner}]
\begin{quote}
\textbf{\pa{ਨੇ}} — \emph{ne} — marks the person who did the action in this frame
\end{quote}

\begin{quote}
\textbf{\pa{ਅਮਨ} \pa{ਨੇ} \pa{ਕੰਮ} \pa{ਕੀਤਾ}} — \emph{Aman ne kamm kiitaa} — Aman did the work
\end{quote}

Put \textbf{\pa{ਨੇ}} immediately after \textbf{\pa{ਅਮਨ}} in this one third-person completed transitive action. The short marker points to the doer. Familiar \textbf{\pa{ਕੀਤਾ}} still agrees with the unmarked masculine-singular object \textbf{\pa{ਕੰਮ}}. Say the whole sentence once, then point to its doer, work, and action.
\end{grammarlens}

\begin{grammarlens}[title={Grammar Lens: source boundary}]
Punjabi University Patiala's \href{https://pt.learnpunjabi.org/lesson14.aspx?l=14}{lesson 14} shows a comparable completed-action pattern and warns that usage varies by sentence. This is a safe small example; not every past sentence takes \textbf{\pa{ਨੇ}}. The same written sound can also be a present-tense auxiliary at the \emph{end} of a different sentence; that is not today's marker. We do not generalise beyond this named-agent, third-person frame.
\end{grammarlens}

\subsection*{Your turn}
\begin{itemize}
\raggedright
  \item \emph{Hear:} \emph{Aman ne kamm kiitaa} Say \textbf{\pa{ਅਮਨ} \pa{ਨੇ} \pa{ਕੰਮ} \pa{ਕੀਤਾ}}.
  \item \emph{Read it:} \textbf{\pa{ਅਮਨ} \pa{ਨੇ} \pa{ਕੰਮ} \pa{ਕੀਤਾ}}
  \item \emph{Write it:} copy only \textbf{\pa{ਨੇ}}, checking \textbf{\pa{ਨ}} then \textbf{\pa{ੇ}}
\end{itemize}

\begin{tcolorbox}[breakable,colback=teal!4,colframe=teal!35!black,title={Before you move on}]
Recall the marker after the named doer.
\end{tcolorbox}
\section[\texorpdfstring{\pa{ਨੇ} (ne)}{ਨੇ (ne)}]{Hear and read one agent marker again}

\label{lesson:PA-C210-ne-review}



\begin{quote}
Remember the tiny marker immediately after \textbf{\pa{ਅਮਨ}} in the completed-action sentence.
\end{quote}

\subsection*{Your turn}
\begin{itemize}
\raggedright
  \item \emph{Hear:} \emph{ne, after the doer} Say \textbf{\pa{ਨੇ}}.
  \item \emph{Read it:} \textbf{\pa{ਨੇ}} Say \emph{ne}.
  \item \emph{Write it:} cover the model and write \textbf{\pa{ਨੇ}} once
\end{itemize}

If it does not come yet, look back at the model and try again. The two visible pieces are familiar; the new work is where this whole marker sits.

\begin{tcolorbox}[breakable,colback=teal!4,colframe=teal!35!black,title={Before you move on}]
Say where the marker belongs: after the named doer in today's completed-action frame.
\end{tcolorbox}
\section[\texorpdfstring{\pa{ਅਮਨ} \pa{ਨੇ} \pa{ਕੰਮ} \pa{ਕੀਤਾ} (Aman ne kamm kiitaa)}{ਅਮਨ ਨੇ ਕੰਮ ਕੀਤਾ (Aman ne kamm kiitaa)}]{Build one completed-action sentence}

\label{lesson:PA-C210-aman-did-work}



\begin{quote}
Say \textbf{\pa{ਅਮਨ}}, then the new marker \textbf{\pa{ਨੇ}}. Recall \textbf{\pa{ਕੰਮ} \pa{ਕੀਤਾ}} from the earlier completed-action chapter.
\end{quote}

\begin{grammarlens}[title={Grammar Lens: bounded frame}]
\begin{quote}
\textbf{\pa{ਅਮਨ} \pa{ਨੇ} \pa{ਕੰਮ} \pa{ਕੀਤਾ}} — \emph{Aman ne kamm kiitaa} — Aman did the work
\end{quote}

Read left to right: named doer, small agent marker, work object, completed action. Keep this one frame; do not use it to guess a rule for every past sentence.
\end{grammarlens}

\subsection*{Your turn}
\begin{itemize}
\raggedright
  \item \emph{Hear:} \emph{Aman did the work} Say \textbf{\pa{ਅਮਨ} \pa{ਨੇ} \pa{ਕੰਮ} \pa{ਕੀਤਾ}}.
  \item \emph{Read it:} \textbf{\pa{ਅਮਨ} \pa{ਨੇ} \pa{ਕੰਮ} \pa{ਕੀਤਾ}} Point to the small marker after \textbf{\pa{ਅਮਨ}}.
  \item \emph{Write it:} copy \textbf{\pa{ਅਮਨ} \pa{ਨੇ} \pa{ਕੰਮ} \pa{ਕੀਤਾ}} once
\end{itemize}

\begin{tcolorbox}[breakable,colback=teal!4,colframe=teal!35!black,title={Before you move on}]
Cover the model, rebuild the four pieces, and check the marker's position.
\end{tcolorbox}
\section[\texorpdfstring{\pa{ਅਮਨ} \pa{ਨੇ} \pa{ਕੰਮ} \pa{ਕੀਤਾ} (Aman ne kamm kiitaa)}{ਅਮਨ ਨੇ ਕੰਮ ਕੀਤਾ (Aman ne kamm kiitaa)}]{A bounded completed-action payoff}

\label{lesson:PA-C210-payoff}



\begin{quote}
Recall the small agent marker, then retrieve the full frame.
\end{quote}

\subsection*{Your turn: four skills}
\begin{itemize}
\raggedright
  \item \emph{Hear:} \emph{Aman ne kamm kiitaa} Say \textbf{\pa{ਅਮਨ} \pa{ਨੇ} \pa{ਕੰਮ} \pa{ਕੀਤਾ}}.
  \item \emph{Read it:} \textbf{\pa{ਅਮਨ} \pa{ਨੇ} \pa{ਕੰਮ} \pa{ਕੀਤਾ}} Give its meaning: Aman did the work.
  \item \emph{Write it:} cover the line and write \textbf{\pa{ਅਮਨ} \pa{ਨੇ} \pa{ਕੰਮ} \pa{ਕੀਤਾ}}
  \item Check that \textbf{\pa{ਨੇ}} is after the named doer, not the final auxiliary.
\end{itemize}

\begin{tcolorbox}[breakable,colback=teal!4,colframe=teal!35!black,title={Before you move on}]
Produce this one third-person completed-action sentence once more. This is not an A2 mock or mock readiness evidence. It does not, by itself, unblock Punjabi A2 mock 1; the full public-corpus closure must be rechecked.
\end{tcolorbox}
